\section{Conformal prediction 用可交换性换有限样本覆盖}
\label{sec:11-Mathematical-Statistics/07-Statistical-Decision-and-Reliable-Prediction/05}

医院运营部门委托你建一个模型，根据入院时的生命体征和化验结果预测患者住院天数。你用 5000 条历史记录训练了一个梯度提升模型，测试 RMSE 2.3 天。现在一个 68 岁的慢阻肺患者入院，模型输出 5.2 天。临床主任问：``有多确定？能给一个区间吗？''

你用线性回归教科书的方法算了一个 \(95\%\) 预测区间：\(5.2 \pm t_{n-p,0.975}\times s\times\sqrt{1+\text{leverage}}\)，得到 \([1.1, 10.7]\)。主任看了看说，这个范围几乎覆盖了科室里绝大部分住院时长，决策价值不大。而且你的梯度提升模型根本不符合高斯误差假设，那个 \(t\) 分位数也不成立。有没有一种方法，不需要假设误差分布，不依赖大样本渐进，也不要求 Bayes 先验，只凭手头这 5000 条数据的某一部分构造出有保的预测区间？

\subsection{Split conformal：拿一部分数据校准，另一部分拟合}

把 5000 条数据拆成两部分：4000 条 training set，1000 条 calibration set。规则很明确——两部分完全隔离。

先在 training data 上拟合任意预测器 \(\widehat f\)。它可以是一个梯度提升模型、一个神经网络、甚至一个人工规则表；conformal 不限制模型形式。

然后在 calibration data 上算 nonconformity scores：

\[
R_i=\abs{Y_i-\widehat f(X_i)},
\qquad i=1,\ldots,m,
\]

其中 \(m=1000\)。这些分数衡量每个 calibration 样本``偏离预测有多远''。把它们从小到大排列：

\[
R_{(1)}\le R_{(2)}\le\cdots\le R_{(m)}.
\]

现在取

\[
k=\left\lceil(m+1)(1-\alpha)\right\rceil.
\]

如果 \(k\le m\)，令 \(q=R_{(k)}\)；如果 \(k=m+1\)，令 \(q=\infty\)——后者只发生在 \(\alpha\) 极小时。对新来的患者 \(x\)，输出预测区间：

\[
C(x)=[\widehat f(x)-q,\;\widehat f(x)+q].
\]

为什么是 \(m+1\) 而不是 \(m\)？因为新样本自己也要参与排名。如果你用 calibration scores 的简单经验 \(1-\alpha\) 分位数，相当于隐含假设新 score 一定不排进前 \(1-\alpha\) 分位之外。把样本量从 \(m\) 修正为 \(m+1\)，正是承认新 score 有同等可能落在任何一个秩位置。取 \(\alpha=0.1\)、\(m=1000\)，则 \(k=\lceil 1001\times 0.9\rceil=901\)，取第 901 小的残差作为 \(q\)。这个 \(m+1\) 修正在一百多个 calibration 点时差别明显，在大样本中影响变小，但它确保了有限样本覆盖的严格性。

\subsection{覆盖保证的论证只依赖一个对称性}

条件于 training data 之后，calibration samples \((X_1,Y_1),\ldots,(X_m,Y_m)\) 和新样本 \((X_{\text{new}},Y_{\text{new}})\) 如果可交换——即这 \(m+1\) 个数据点的联合分布在任意排列下不变——那么它们的 nonconformity scores \(R_1,\ldots,R_m,R_{\text{new}}\) 也可交换。

可交换性的直接后果是：新 score 的秩在 \(1,2,\ldots,m+1\) 上均匀分布。因此

\[
P\bigl(R_{\text{new}}\text{ 的秩 }\le k\bigr)
=\frac{k}{m+1}
\ge 1-\alpha.
\]

\(R_{\text{new}}\) 的秩 \(\le k\) 当且仅当 \(R_{\text{new}}\le R_{(k)}=q\)，即 \(\abs{Y_{\text{new}}-\widehat f(X_{\text{new}})}\le q\)，也就是 \(Y_{\text{new}}\in C(X_{\text{new}})\)。所以

\[
P\bigl(Y_{\text{new}}\in C(X_{\text{new}})\bigr)\ge 1-\alpha.
\]

这个推导中没有任何关于模型形式的假设，没有渐近符号，没有高斯分布，没有先验。它只要求可交换性——比独立同分布稍弱（i.i.d. 的序列可交换，但可交换的序列不一定独立），却足以把秩的均匀性变成覆盖保证。有 ties 时采用保守随机化处理，保证依然成立。

注意保证给了什么、没给什么：``模型越好，区间越窄''——如果预测误差普遍小，第 \(k\) 分位数就小，区间就紧。但即使模型糟糕到残差巨大，覆盖保证依然成立，只不过区间宽到丧失临床意义。Conformal 校准的是覆盖概率，不是预测精度。

\subsection{分类问题输出集合而非单标签}

分类器同样可以用 conformal。设模型对每个类别输出概率估计 \(\widehat p_y(x)\)，定义 nonconformity 为模型给真实类别概率的补数：

\[
R_i=1-\widehat p_{Y_i}(X_i).
\]

分数小的样本是``模型很有把握且判断正确''的样本；分数大的样本则是模型对真实类别信心不足甚至完全看错。用 calibration quantile \(q\) 构造预测集合：

\[
C(x)=\set{\,y:1-\widehat p_y(x)\le q\,}.
\]

对容易的样本，\(1-\widehat p_{\text{true}}(x)\) 很小，\(q\) 过滤后可能只留一个类别；对模棱两可的样本，多个类别的 \(\widehat p_y(x)\) 都差不多，集合中保留下来的类别变多。极端情况下——模型对所有类别都没信心或 calibration 分数差异极大——集合可能为空，也可能包含全部类别。这并非缺陷，而是说明模型在当前样本上能提供的分辨信息很少。集合大小本身就是不确定性的一种度量。

更精细的构造可以让区间或集合适应异方差和类别不平衡：用归一化残差 \(\abs{Y_i-\widehat f(X_i)}/\widehat\sigma(X_i)\) 作为 score 能让区间在预测不确定的区域自动变宽；对不同类别分别校准可以防止多数类的分数主导了 quantile 的选择。但每次修改 score 函数，都要重新说明可交换性单位：校准集和新样本的 scores 是否仍然可交换？归一化残差中的 \(\widehat\sigma\) 是用哪些数据拟合的？如果它依赖 calibration data 本身，对称性就变了。

\subsection{边际覆盖不是每个子群的覆盖}

覆盖保证中的概率是平均于新样本的整体分布：

\[
P\bigl(Y_{\text{new}}\in C(X_{\text{new}})\bigr)\ge 1-\alpha.
\]

它不能推出对每个固定的 \(x\) 都有

\[
P\bigl(Y\in C(X)\mid X=x\bigr)\ge 1-\alpha.
\]

更不能保证每个年龄组、每种合并症组合、每家医院都达到 \(90\%\) 的名义覆盖率。一个可能发生的情形是：年轻低风险患者群中覆盖高达 \(97\%\)，高龄多病患群中只有 \(78\%\)，二者均值仍有 \(90\%\)。整体覆盖被容易的群体补偿了困难群体。

在无分布假设下要求所有点都达到精确的条件覆盖是不可能的任务。局部条件覆盖需要额外的平滑性假设、足够的组内样本量或模型结构信息。因此 conformal 结果的使用者应预先指定有实际意义的子群——不是看完结果后无限搜索最差子群——分别检查它们的覆盖率、集合宽度与覆盖的不确定性。

\subsection{分布漂移和数据复用会摧毁论证}

可交换性一旦被打破，整条秩论证链条就断了。如果 calibration data 全部来自 2023 年的 A 医院，而新预测要用在 2025 年换了检验设备的 B 医院，scores 不再可交换，名义 \(90\%\) 覆盖可能掉到 \(60\%\) 以下，而且你不知道掉了多少。

某些 covariate shift 场景可以用加权 conformal 缓解：给 calibration samples 加权来逼近目标分布的 score 分布。但这需要可靠的密度比估计，并且要求目标分布的支持包含在 calibration 分布的支持内——如果 B 医院设备测出的某项指标范围完全超出 A 医院的历史范围，加权也无济于事。Concept drift（条件分布 \(Y|X\) 本身改变）更不能用简单权重修复。

同样的逻辑也适用于数据复用。如果同一个 calibration set 被你反复用来选择模型架构、挑选 score 函数、调试超参数，然后才决定最终的 conformal quantile，那 calibration set 实际上参与了全部建模选择，不再是``只用于校准的隔离数据''。新样本的 rank 不再与一个预先固定下来的程序直接对称。正确做法是把数据分成三段——训练集用于拟合模型，校准集用于确定 conformal quantile，测试集用于最终验证——或者使用能够还原完整选择流程的方案（如 cross-conformal 或 jackknife+），但那些方案会引入各自的有限样本行为和假设。

\subsection{能相信什么，需要额外验证什么}

Conformal prediction 提供的是一个有限样本的分布无关覆盖保证，代价是放弃了条件覆盖和模型改进。你得到的保证是你的预测区间至少以概率 \(1-\alpha\) 包含新观测——不多不少。以下判断不能由 conformal 直接推出，而需要额外的验证：

\begin{itemize}
\tightlist
\item
  每个预设子群的覆盖是否达标——需要在子群层面单独统计。
\item
  区间的平均宽度是否窄到有临床决策价值——覆盖和宽度是一对需要同时审视的指标，单报覆盖而隐藏平均宽度等于隐瞒了方法付出的信息代价。
\item
  分布随时间漂移后覆盖是否仍然成立——需要在新时期的新数据上用相同的校准程序重新计算实际覆盖，而不是依赖旧日的保证。
\end{itemize}

Conformal 最大的诚实在于：它把全部假设集中到一个明确的条件上——可交换性——然后从这一个条件严格推出有限样本覆盖。你认为数据满足可交换性，覆盖保证就成立；你认为不满足，就失去了保护。它不是``无需假设''的魔法，而是把假设显式地放在了你能审查的地方。
